\section{A global game to understand stablecoin risk dynamics}\label{sec:global_game_policy}

A stablecoin holder's decision to convert into the underlying fiat is not an isolated risk-management problem. Par convertibility fails through two channels, and a holder who anticipates others converting at the same time faces a worse secondary price, so the decision is a coordination problem under imperfect information rather than an individual one. We model it as a global game in the spirit of \citet{morris_shin_1998}, adapted to the dual-channel exit through the issuer's primary redemption and a secondary market with depth-dependent price impact that distinguishes a stablecoin from a deposit. Figure~\ref{fig:global_game} shows the coordination structure; the formal model is in Appendix~\ref{app:formal_model}.

\begin{figure}[p]
\centering
\includegraphics[width=\textwidth]{figures/fig_morris_shin_clock_preview.pdf}
\caption{Coordination structure of the stablecoin redemption problem. Holders face two distinct risks: counterparty risk on the issuer's reserve and depeg risk on the secondary conversion market. The depeg is a coordination outcome: conversions exhaust the issuer's primary redemption and depress the secondary price, a self-fulfilling run. Each holder observes the same fundamental $\theta$ with an idiosyncratic noisy signal $s_i$, picks a redemption threshold $s^{*}$ given the aggregate run she expects from the rest, while the aggregate run itself is determined by the threshold every holder uses. Equilibrium is the unique fixed point $s^{*}$ such that $B(s^{*}) = s^{*}$. The formal derivation is in Appendix~\ref{app:formal_model}.}\label{fig:global_game}
\end{figure}

\subsection{Why a global game}\label{sub:why_different}

Each comparison in the corridor exercise pits a stablecoin against a non-stablecoin alternative for the same payment task: a public instant payment system, a card network, a fintech remittance service, or a correspondent-bank wire. None is operationally risk free, but all share one property: the value sent and the value accepted are measured in the unit of account, with no exchange-rate-like discount embedded in the channel. A stablecoin is different. The holder carries a private-issuer liability whose promise of par convertibility can fail in two ways no other channel here exhibits: the issuer may lack the liquid reserves to redeem at par (counterparty risk), and the secondary venue where unmet redemptions clear has finite depth, so the realised price falls below par when conversion demand is high (depeg risk). The depeg channel is itself a coordination problem: conversions exhaust the issuer's primary redemption and depress the secondary price, a self-fulfilling run. Counterparty risk is a distinct issuer-level component, and at the calibration the two price additively, so their interaction is second-order. \citet{aldasoro_mehrling_neilson_2023} frame the same object from a money-view perspective, and \citet{anadu_etal_2023_runs} document the run dynamics empirically, including the break-the-buck threshold near par.

The implication for a central bank is direct. A comparison that prices the stablecoin at its transfer fee plus a static risk premium misstates its position in the cost ranking: too low a premium makes the stablecoin look attractive when it carries structural tail risk a public rail does not, and too high a premium rules it out by assumption. Neither informs the decision a central bank actually faces over its payment infrastructure: what to build, who to license, and how to supervise.

Pricing the depeg risk requires modelling the redemption decision as a coordination problem. A holder of one unit decides whether to convert at the issuer or hold the unit and use it for payment. The value of holding depends on the latent quality of the issuer's reserve, governance, and standing, collected into a single fundamental $\theta$ that holders observe only through noisy private signals $s_i = \theta + \sigma \varepsilon_i$. The value of converting depends on what others do: if many convert at once, liquid reserves are exhausted, the residual demand spills into the secondary market, and the realised conversion price falls. The marginal holder's payoff therefore depends on the aggregate conversion mass, which itself depends on every holder's signal about the same $\theta$. \citet{morris_shin_1998} show that this class of problem has a unique equilibrium in threshold strategies: every holder converts if and only if her signal falls below a critical value $s^*$, set endogenously where the marginal holder is indifferent given that all others use the same rule. The equilibrium is the fixed point at which the threshold is consistent with the conversion mass it implies; we solve it numerically (Section~\ref{sec:formal_model_summary} and Appendix~\ref{app:formal_model}).

Four features make the global game the natural model rather than a stylised analogy: the action is binary and the payoff to converting falls as others convert, the strategic complementarity that delivers a unique threshold; information is heterogeneous across market makers, retail, and cross-border holders; the conversion venue is a secondary market of limited, stress-contracting depth whose elasticities are identified from the March 2023 price dynamics; and the secondary price and the conversion mass are jointly determined at the fixed point, so that the total risk decomposes into a counterparty component, a baseline-regime depeg component, and a stress-regime depeg component. Throughout, holders earn zero nominal yield, consistent with the GENIUS Act and MiCA prohibitions on interest on payment stablecoins; the only return to holding is the convenience yield $\omega$ of payment use, formalised in Appendix~\ref{app:formal_model}.\footnote{The Coinbase ``USDC Rewards'' program, paid out of reserve revenue rather than as issuer interest, is being narrowed toward transaction-linked rewards under the GENIUS Act and parallel OCC guidance. The zero-yield assumption is consistent with the prevailing regulatory baseline; relaxing it as a future-state extension would amplify, rather than dilute, the depeg-channel results derived below.}

\subsection{Stablecoins are not a homogeneous asset class}\label{sub:not_homogeneous}

A further observation matters for the empirical comparison: the two largest fiat-backed stablecoins (USDC and USDT, jointly carrying over 90\% of the segment's market value) have structurally different primitives. The differences are not stylistic; they map directly to different parameter values in the global game and therefore to different risk premiums.

\textbf{Reserve composition.} USDC is backed exclusively by short-dated US Treasuries and cash deposits at G-SIBs, with monthly attestations and daily disclosure of holdings since 2022. USDT holds a heterogeneous mix of US Treasuries, secured loans, and historically commercial paper; attestations are produced by BDO without the granularity of a full audit. The implication for the model is that USDC's loss-given-default parameter $\xi$ is structurally lower than USDT's, and USDT's hazard rate $h$ is structurally higher, reflecting both the historical enforcement record (NYAG and CFTC settlements in 2021) and the residual uncertainty in attestation-only disclosure.

\textbf{Primary market structure.} USDC offers permissionless primary redemption via Circle Account with daily settlement and immaterial minimums; USDT requires whitelisting, applies a USD 100,000 minimum on direct redemption, and operates on a slower cycle. The implication is that USDC's primary redemption capacity $R$ is larger and its signal noise $\sigma$ is higher because the population of holders is more atomistic and less informed; USDT's $R$ is smaller and its $\sigma$ is structurally lower because arbitrage activity concentrates in a few large market makers (Cumberland, Wintermute, B2C2).

\textbf{Secondary market depth.} USDT is the most liquid stablecoin globally, with primary venue Binance and per-side depth typically above USD 100 million; USDC's primary venue is Coinbase with depth in the USD 50 million range. The implication is that USDT's price-impact scale $\psi$ is structurally smaller, its depth-erosion parameter $\zeta$ is smaller, and the structural backstop $A$ provided by concentrated arbitrage is larger.

\textbf{Idiosyncratic exposure.} USDC's March 2023 depeg was driven by direct exposure to a bank failure (SVB held a portion of the cash leg of Circle's reserve). USDT's May 2022 depeg was driven by contagion from the Terra/UST collapse and the subsequent reputational pressure on every dollar stablecoin, not by Tether-specific reserve impairment. The two stress episodes are structurally different and live in different regimes of the model.

The empirical comparison in Section~\ref{sec:empirical_comparison} accordingly applies a USDC-calibrated risk premium to use cases where USDC is the natural stablecoin choice (EU-originated cross-border under MiCA, institutional settlement) and a USDT-calibrated premium to use cases where USDT dominates the real-world flow (LATAM remittance and domestic retail). The implied per-asset risk profiles are inverted in composition: USDC is depeg-dominant (a clean reserve carrying a residual risk of secondary-market dispersion when convertibility is briefly questioned), while USDT is counterparty-dominant (a structurally larger residual hazard rate, partially offset by a deeper and more concentrated secondary market that limits depeg propagation). The policy implication is that "stablecoins" as a category is too coarse: the regulator needs to distinguish the structural primitives of each issuer.

\subsection{An empirical comparison}\label{sub:what_it_delivers}

The output of the global-game stage, for each asset under its own calibration, is three numbers in basis points: a counterparty-loss component $\ell_C$, a depeg-loss baseline component $\ell_{D,b}$, and a depeg-loss stress component $\ell_{D,s}$. Each is an expected per-unit loss conditional on the regime. Their probability-weighted sum is the total expected risk premium the holder bears for using the stablecoin as a payment vehicle:
\begin{equation}\label{eq:total_risk_premium}
\Pi_{\text{risk}} = \ell_C + (1-p_s)\ell_{D,b} + p_s \ell_{D,s},
\end{equation}
where $p_s$ is the probability of the stress regime, calibrated to historical episode frequency. The two stablecoins receive separate calibrations and deliver premiums with inverted compositions; Section~\ref{sec:identification} reports the magnitudes and Sections~\ref{sec:empirical_comparison} and~\ref{sec:discussion} carry them through to the corridor comparison and its supervisory implication.

This premium enters the empirical comparison in Section~\ref{sec:empirical_comparison} as the stablecoin's risk loading. The competitor rails carry no counterparty or depeg premium, only the operational risk common to every payment system, which washes out across the comparison. The stablecoin's total cost adds to the on-chain transfer fee the off-ramp basis in the destination jurisdiction and the structural premium $\Pi_{\text{risk}}$; a competitor rail adds instead the working-capital cost of its settlement delay. Section~\ref{sec:empirical_comparison} writes the two cost expressions out and applies them corridor by corridor.

The risk premium $\Pi_{\text{risk}}$ does not vanish under any operational improvement in the stablecoin transfer technology. It reflects the economic substance of holding a private-issuer liability rather than a sovereign-backed unit of account, and an analyst comparing nominal transfer fees alone would not observe it. The model quantifies the premium consistent with the price dynamics observed in the calibration episode.

The premium $\Pi_{\text{risk}}$ prices counterparty and depeg risk structurally. Future regulatory change enters the comparison through a separate channel: it reprices the observable off-ramp basis on each leg, not the structural premium. \citet{copestake_etal_2026} document that policy enactments reprice the equity value of incumbent payment firms, and the descriptive evidence on the two most recent stablecoin-side packages (MiCA Title V and Banco Central do Brasil Resolution 561/2026), reported in the internet appendix, shows no measurable basis response, with incumbent rates adjusting only at a competitive lag. The framework prices regulatory change once it registers in the prices it observes.

\subsection{Formal model and Propositions}\label{sec:formal_model_summary}

The formal model occupies Appendix~\ref{app:formal_model}; we state here the objects needed to read the results. The model has four layers: a payment-rail block, in which the issuer redeems the first $R$ units of demand at par out of liquid reserves; an issuer balance-sheet layer carrying the reserve quality, the default hazard $h$, and the loss given default $\xi$; a secondary-market layer on which unmet demand clears at a price that falls, convexly, in the volume of failed redemptions and in a depth that itself contracts under stress; and a coordination layer that ties the three together.

The information state is a single private signal. A continuum of holders share the latent issuer fundamental $\theta$, and each observes
\[
s_i = \theta + \sigma\,\varepsilon_i, \qquad \varepsilon_i \sim \mathcal{N}(0,1),
\]
so the posterior is $\theta \mid s_i \sim \mathcal{N}(s_i, \sigma^2)$ and the signal noise $\sigma$ is the only source of imperfect information. Holders play threshold strategies, converting when $s_i < s^*$, so the mass that attempts redemption at fundamental $\theta$ is
\[
D(\theta; s^*) = \Phi\!\left(\frac{s^* - \theta}{\sigma}\right),
\]
with $\Phi$ the standard normal CDF. Writing $V_R(D)$ for the payoff to converting (par on the first $R$ units, the depressed secondary price on the rest) and $V_H(\theta)$ for the payoff to holding (the fundamental plus the convenience yield $\omega$), the equilibrium threshold is the endogenous fixed point at which the marginal holder is indifferent,
\[
F(s^*) \;:=\; \mathbb{E}_{\theta \mid s_i = s^*}\!\big[\, V_R(D(\theta; s^*)) - V_H(\theta) \,\big] \;=\; 0.
\]
The threshold $s^*$ enters $F$ both as the conditioning event of the posterior and through the aggregate mass $D$ it induces; this two-way dependence is what makes the fixed point endogenous, and it admits no closed form. This is where the model departs from the rail-side global game of \citet{klee_etal_2026_fragility}, whose congestion threshold is available in closed form under linear congestion costs and a uniform signal: here the convex, stress-contracting secondary depth and the endogenous coordination fixed point preclude a closed form, and the premium is recovered numerically. The explicit secondary-price formation and the payoffs $V_R$ and $V_H$ are in Appendix~\ref{app:formal_model}.

Three propositions, stated and proved in Appendix~\ref{app:formal_model}, give the model the properties the calibration relies on. Proposition~\ref{prop:existence} establishes that the redemption threshold is unique: under a noisy signal and a convex price-impact schedule, the best-response map that sends a conjectured threshold to the holder's optimal reply is a contraction, so it has exactly one fixed point. Uniqueness is what makes the exercise well posed. A model with multiple equilibria would deliver a range of redemption intensities for the same fundamentals and could not map an observed price path into a single risk premium; the contraction property, which we confirm numerically at the calibrated parameters, rules that out and lets the calibration read one premium off the data.

The other two propositions describe how the equilibrium responds to the primitives that distinguish one issuer from another, and they connect the model to the empirical literature on stablecoin runs. Proposition~\ref{prop:reserves} signs the response to reserve strength: a deeper reserve sustains a higher secondary price for any given conversion volume, which raises the value of converting and so weakly raises the equilibrium conversion mass, in the spirit of the reserve-quality mechanism in \citet{ma_zeng_zhang_2025}. Proposition~\ref{prop:convenience} signs the response to the payment usefulness of the coin: a higher convenience yield raises the value of continuing to hold, contracts the set of fundamentals at which the marginal holder prefers to convert, and so lowers the equilibrium conversion mass, the adoption-side counterpart of the fragility channel in \citet{bertsch_2025_adoption_fragility}. The two comparative statics are the formal content of the claim in Section~\ref{sub:not_homogeneous} that reserve composition and payment use are not stylistic differences between USDC and USDT but primitives that move the equilibrium and therefore the priced premium.

The three risk components that enter the empirical comparison are expectations taken at the solved equilibrium, and the formal derivations guarantee that they are well-defined functions of the moments of the observed price process rather than free parameters.
